\documentclass{szb-solution}

\area{Integrálszámítás}
\title{Integrálszámítás III}
\subject{Matematika G1}
\subjectCode{BMETE94BG01}
\date{Utoljára frissítve: \today}
\docno{12}

\begin{document}
\maketitle

\subsection{Trigonometrikus integrálok}

Minden határozatlan integrálban $C\in\mathbb R$ tetszőleges állandó.

\begin{enumerate}[a)]
  \item Az $u=\cos x$, $\dd u=-\sin x\dd x$ helyettesítéssel
        $$
          \boxed{\int\cos^3x\sin x\dd x=-\frac{\cos^4x}{4}+C}
          \text.
        $$

  \item Egy koszinuszt leválasztunk, a többit szinusz segítségével írjuk fel:
        $$
          \cos^5x=(1-\sin^2x)^2\cos x
          \text.
        $$
        Az $u=\sin x$ helyettesítés után
        $$
          \boxed{\int\cos^5x\dd x
          =\sin x-\frac23\sin^3x+\frac15\sin^5x+C}
          \text.
        $$

  \item A hatványcsökkentő azonosságokkal
        $$
          \sin^4x\cos^2x
          =\frac1{16}-\frac1{32}\cos2x
           -\frac1{16}\cos4x+\frac1{32}\cos6x
          \text.
        $$
        Így
        $$
          \boxed{\int\sin^4x\cos^2x\dd x
          =\frac{x}{16}-\frac{\sin2x}{64}
           -\frac{\sin4x}{64}+\frac{\sin6x}{192}+C}
          \text.
        $$
\end{enumerate}

\subsection{Összetett integrálási feladatok}

\begin{enumerate}[a)]
  \item Legyen $t=\sqrt x$, így $x=t^2$ és $\dd x=2t\dd t$. Parciális
        integrálással
        $$
          2\int t\sin t\dd t=-2t\cos t+2\sin t+C
          \text,
        $$
        ezért
        $$
          \boxed{\int\sin\sqrt x\dd x
          =-2\sqrt x\cos\sqrt x+2\sin\sqrt x+C}
          \text.
        $$

  \item Legyen $t=\ln x$, ekkor $\dd t=\dd x/x$. A valós logaritmusok miatt
        $x>1$, és
        $$
          \boxed{\int\frac{\ln\ln x}{x}\dd x
          =\ln x\,\ln\ln x-\ln x+C}
          \text.
        $$

  \item Az abszolútérték miatt intervallumonként integrálhatunk:
        $$
          \int|x|\dd x=
          \begin{cases}
            -x^2/2+C,&x<0,\\
            x^2/2+C,&x\geq0.
          \end{cases}
        $$
        Ez egyetlen képlettel
        $$
          \boxed{\int|x|\dd x=\frac{x|x|}{2}+C}
          \text.
        $$

  \item Legyen $t=x^x$. Mivel $x>0$ esetén
        $$
          \dd t=x^x(\ln x+1)\dd x=t(\ln x+1)\dd x,
        $$
        ezért
        \begin{align*}
          \int\frac{\ln x+1}{x^x-1}\dd x
          &=\int\frac{\dd t}{t(t-1)}
            =\int\left(-\frac1t+\frac1{t-1}\right)\dd t\\
          &=\boxed{\ln\left|\frac{x^x-1}{x^x}\right|+C}
          \text.
        \end{align*}

  \item Legyen $t=2x-1$, így $x=(t+1)/2$ és $\dd x=\dd t/2$. Ekkor
        $$
          x^2-3x+2=\frac{t^2-4t+3}{4}
          \text,
        $$
        tehát
        $$
          \int(x^2-3x+2)\sqrt{2x-1}\dd x
          =\frac18\int(t^{5/2}-4t^{3/2}+3t^{1/2})\dd t
          \text.
        $$
        Visszahelyettesítve
        $$
          \boxed{\frac{(2x-1)^{7/2}}{28}
          -\frac{(2x-1)^{5/2}}5
          +\frac{(2x-1)^{3/2}}4+C}
          \text.
        $$
\end{enumerate}

\subsection{Határozott integrálok}

\begin{enumerate}[a)]
  \item Egy teljes perióduson
        $$
          \boxed{\int_0^{2\pi}\cos x\dd x
          =[\sin x]_0^{2\pi}=0}
          \text.
        $$

  \item Parciális integrálással
        $$
          \int_0^1x\sinh x\dd x
          =[x\cosh x]_0^1-\int_0^1\cosh x\dd x
          =\cosh1-\sinh1
          \text.
        $$
        Mivel $\cosh1-\sinh1=e^{-1}$,
        $$
          \boxed{\int_0^1x\sinh x\dd x=\frac1e}
          \text.
        $$

  \item Az $y=\sqrt{9-x^2}$ görbe a $3$ sugarú kör felső félköre, ezért
        $$
          \boxed{\int_{-3}^3\sqrt{9-x^2}\dd x=\frac{9\pi}{2}}
          \text.
        $$
        Ugyanez az $x=3\sin t$ helyettesítéssel is megkapható.
\end{enumerate}

\subsection{Polinom és az x-tengely által bezárt terület}

Legyen
$$
  f(x)=(x+1)x(x-2)=x^3-x^2-2x
  \text.
$$
A zérushelyek $-1$, $0$ és $2$, így a két korlátos tartomány a $[-1,0]$ és
$[0,2]$ intervallumokon található. Egy primitív függvény
$$
  F(x)=\frac{x^4}{4}-\frac{x^3}{3}-x^2
  \text.
$$
Az első intervallumon $f\geq0$, a másodikon $f\leq0$, ezért
\begin{align*}
  T
  &=\int_{-1}^0f(x)\dd x-\int_0^2f(x)\dd x\\
  &=\bigl(F(0)-F(-1)\bigr)-\bigl(F(2)-F(0)\bigr)\\
  &=\frac5{12}+\frac83
  =\boxed{\frac{37}{12}}
  \text.
\end{align*}

\subsection{Két görbe által bezárt terület}

Az $f(x)=x^4$ és $g(x)=3x^2-2$ görbék metszéspontjainak abszcisszái:
$$
  x^4=3x^2-2
  \quad\Longleftrightarrow\quad
  (x^2-1)(x^2-2)=0
  \quad\Longleftrightarrow\quad
  x=\pm1,\ \pm\sqrt2
  \text.
$$
Mindkét függvény páros. A $[0,1]$ intervallumon $f\geq g$, az
$[1,\sqrt2]$ intervallumon $g\geq f$. A három korlátos tartomány teljes
területe tehát
\begin{align*}
  T
  &=2\left(
    \int_0^1(f-g)\dd x+
    \int_1^{\sqrt2}(g-f)\dd x\right)\\
  &=2\left(\frac65+\frac65-\frac{4\sqrt2}{5}\right)\\
  &=\boxed{\frac{24-8\sqrt2}{5}}
  \text.
\end{align*}

\subsection{A kör görbevonalú trapézének területe}

A paraméterezés
$$
  x(t)=a\cos t,
  \qquad
  y(t)=a\sin t,
  \qquad
  0\leq t\leq2\pi
  \text,
$$
ahol $\dot x(t)=-a\sin t$. A görbevonalú trapéz területe
\begin{align*}
  T
  &=\int_0^{2\pi}|\dot x(t)y(t)|\dd t
    =a^2\int_0^{2\pi}\sin^2t\dd t\\
  &=a^2\int_0^{2\pi}\frac{1-\cos2t}{2}\dd t
    =\boxed{\pi a^2}
  \text.
\end{align*}
Ez természetesen megegyezik az $a$ sugarú kör területével.

\end{document}
